%/*% pdflatex -file-line-error -interaction nonstopmode %

\documentclass{article}
\usepackage{amsmath}
\usepackage{amsfonts}
\usepackage{hyperref}
\usepackage{exercises}
\usepackage{amsthm}
\usepackage{yhmath}
\newcommand*{\oldneg}{\mathord{\sim}}

\theoremstyle{definition}
\newtheorem{definition}{Definition}

\setlength{\parindent}{0pt}
\setlength{\parskip}{10pt}

\title{1.2 Exercises}
\begin{document}
\maketitle


\section*{1}
Vertex form.

\section*{2}
In intercept form, we are given $p$ and $q$, the two $x$-intercepts of the parabola. By using the definition of the axis of symmetry, we can conclude 
that the $x$-coordinate of the vertex is $\frac{p+q}{2}$. Then, we can plug this value for $x$ in the equation we are given, thus yielding the $y$-coordinate
 of the vertex.

\section*{12}
B

\section*{22}
First, we find the $x$-coordinate of the vertex, which is the average of the two $y$-intercepts, that is: $\frac{6-4}{2}=1$. Then, the $x$-coordinate of the 
vertex is $-(1-6)(1+4)=25$. Thus, the answer is A.

\section*{23}
The intercept form of a parabola is $y=a(x-p)(x-q)$, which tells us that, in the given equation, the two intercepts are actually $2$ and $-3$.

\section*{24}
We can solve this by multiplying the two binomials: \newline $y=(x+4)(x+3)\newline\rightarrow =x^2+3x+4x+12\newline\rightarrow =x^2+7x+12$

\section*{25}
We can solve this like the above: \newline $y=(x-5)(x+3)\newline\rightarrow =x^2+3x-5x-15\newline\rightarrow =x^2-2x-15$
\section*{26}
We can solve this like the above problem: \newline $y=4(x+1)(x-6)\newline\rightarrow =4(x^2-6x+x-6)\newline\rightarrow =4x^2-20x-24$

\section*{30}
We can solve this by rewriting the squared binomial, and then multiplying out all of the terms:\newline
$g(x)=-(x+6)^2+10\newline\rightarrow =-(x+6)(x+6)+10\newline\rightarrow =-(x^2+6x+6x+36)+10\newline\rightarrow =-x^2-12x-26$

\section*{31}
$y=5(x+3)^2-4\newline\rightarrow =5(x+3)(x+3)-4\newline\rightarrow =5(x^2+3x+3x+9)-4\newline\rightarrow =5x^2+30x+41$

\section*{32}
$f(x)=12(x-1)^2+4\newline\rightarrow =12(x-1)(x-1)+4\newline\rightarrow =12(x^2-x-x+1)+4\newline\rightarrow =12x^2-24x+16$

\section*{38}
Because the value of $a$ is negative, this function has a maximum value. First, the $x$-coordinate of the vertex is $\frac{0+9}{2}=4.5$, and thus, 
the $y$-coordinate of the vertex, and the maximum value, is $-12(4.5)(4.5-9)=243$

\section*{39}
Because the value of $a$ is positive, this function has a minimum value. First, the $x$-coordinate of the vertex is $\frac{0-15}{2}=-7.5$, and thus, 
the $y$-coordinate of the vertex, and the minimum value, is $8(-7.5)(-7.5+15)=-450$

\section*{40}
Because the value of $a$ is positive, this function has a minimum value. First, the $x$-coordinate of the vertex is $\frac{3+6}{2}=4.5$, and thus, 
the $y$-coordinate of the vertex, and the minimum value, is $2(4.5-3)(4.5-6)=-4.5$

\section*{41}
Because the value of $a$ is negative, this function has a maximum value. First, the $x$-coordinate of the vertex is $\frac{-9+4}{2}=-2.5$, and thus, 
the $y$-coordinate of the vertex, and the maximum value, is $-5(-2.5+9)(-2.5-4)=211.25$

\section*{50} 
Vertex form in standard form: $y=a(x-h)^2+k$\newline$\rightarrow =a(x-h)(x-h)+k$\newline$\rightarrow =a(x^2-xh-xh+h^2)+k$
\newline$\rightarrow =ax^2-2ahx+ah^2+k$
\newline 
Hence, the $x$-coordinate of the vertex of this function occurs at the line $x=-\frac{-2ah}{2a}=h$, and thus, the $y$-coordinate of the vertex of this 
function is $ah^2-2ahh+ah^2+k=k$, meaning the vertex of this function is $(h,k)$.\newline
Intercept form in standard form: $y=a(x-p)(x-q)\newline\rightarrow =a(x^2-px-qx+pq)\newline\rightarrow =ax^2-apx-aqx+apq\newline\rightarrow 
=ax^2-(ap+aq)x+apq\newline$ Hence, the $x$-coordinate of the vertex of this function is $x=-\frac{-(ap+aq)}{2a}=\frac{p+q}{2}$, and thus, the vertex of 
this function occurs on the line $x=\frac{p+q}{2}$

\section*{51}
We are given a function with a negative value for $a$, and hence it has a maximum value, which is the $y$-coordinate of the vertex. The vertex is 
$(14,6)$, and hence, the maximum height that the kangaroo can jump is $6$ ft. Because we know that the vertex's $x$-coordinate would be halfway 
between where the kangaroo jumps, and where it lands, we can conclude that the kangaroo's jump is $28$ ft. wide.

\section*{52}
The vertex of the given function is $(52.5,45)$. Because we know that the vertex is $52.5$ meters from the left end of the arch, the axis of symmetry lets 
us conclude that it is $52.5$ meters from the right end of the arch, and thus, that the width of the arch is $105$ meters.

\section*{53}
We are given a quadratic equation in intercept form, with the two intercepts being $0$ and $160$. Hence, the field's width is $160$ ft. Also, the $x$
-coordinate of the vertex is $\frac{160}{2}=80$, and thus, the $y$-coordinate, and the maximum height of the field, is $-0.000234(80)(80-160)=1.4976$ ft.

\section*{54}
The vertex of the first equation is $(6,18)$, and thus it's maximum height is $18$. The vertex of the second equation is $(6,42)$, and hence it's 
maximum height is $42$. Thus, the maximum height of the pogo stick with a bow spring is the larger of the two. Because these two equations are written in 
vertex form, we are given that the vertex is $(h,k)$ in their respective equations. Hence, changing the constant $a$ will not alter the vertex, and hence 
will not alter the maximum height. Also, because the maximum height is defined to be $k$, chaning $h$ will not affect the maximum height. Thus, the only variable 
that affects the value of the maximum height is $k$.

\section*{55}
(a): We can tell that the $x$-coordinate of the vertex in this equation is $\frac{5.52+22.6}{2}=14.06$, and that the $y$-coordinate of the vertex is 
$-0.761(14.06-5.52)(14.06-22.6)\approx 55.5$\newline Hence, a moisture of $14.06$ contains the maximum popping volume of $55.5\frac{cm^3}{g}$\newline
(b): We can tell that the $x$-coordinate of the vertex in this equation is $\frac{5.35+21.8}{2}=13.575$, and that the $y$-coordiante of the vertex is 
$-0.652(13.575-5.35)(13.575-21.8)\approx 44.1$. Hence, a moisutre of $13.575$ contains the maximum popping volume of $44.1\frac{cm^3}{g}$\newline
(c): For the first function, the domain is $5.52\le x\le 22.6$. This is because, any other $x$ values yield negative values, which do not make sense 
in this situation. The range is $0\le y\le 55.5$, as $55.5$ is the maximum value. We can use similar reasoning for the second equation, to yield a 
domain of $5.35\le x\le 21.8$ and a range of $0\le y\le 44.1$

\section*{56}
We are given the vertex, as well as one point, namely $(0,0)$. So, we can solve for $a$:\newline
$0=a(0-33)^2+5\newline\rightarrow a\approx -0.0046$\newline
So, the equation to model this situation is $y=-0.0046(x-33)^2+5$ Chaning the value of $a$ would either make the parabola wider, or more narrow. Changing 
$h$ by $n$ units would shift the parabola to the right by $n$ units. Changing $k$ by $m$ units would shift the parabola up by $m$ units.
\end{document}

